% ============================================================================
% 02_architecture | 05_normalization
% Section-body file: include immediately after the relevant \section command
% in the master tstextbook.tex file.
%
% This file intentionally contains no document class, package imports,
% \chapter, \section, or bibliography command. Figure source fragments are
% stored in the root-level figures directory and inserted with \input.
% ============================================================================

The preceding discussion established how feed-forward transformations and
activation functions modify the feature representation of each token. These
operations are repeatedly composed with attention sublayers throughout a
Transformer. Their repeated application raises a further architectural
question: how should the numerical scale of the intermediate representations
be controlled as information passes through a deep sequence of transformations?
Normalization addresses this question by transforming activations according
to statistical quantities calculated from selected dimensions of the
representation. Its position within Transformer blocks, its mathematical
properties, and its alternatives are therefore necessary parts of the
architecture before the attention mechanism is considered in full.

\subsection{Architectural motivation and statistical foundations}
\label{subsec:norm-motivation-statistics}

A neural network transforms its input through a sequence of parameterized
operations. In a Transformer, the representations passed between these
operations are vectors associated with individual tokens. Each vector contains
features that encode information accumulated by the model at a particular
stage of computation. The numerical values of these features are called
activations. Their magnitudes and distributions affect the inputs received by
subsequent transformations and the derivatives through which the parameters
are optimized.

Consider a batch of token sequences represented by the activation tensor
\begin{equation}
    \mathbf{X}\in\mathbb{R}^{B\times T\times C},
    \label{eq:norm-input-tensor}
\end{equation}
where \(B\) denotes the batch size, \(T\) the number of token positions in
each sequence, and \(C\) the hidden or embedding dimension. The scalar
activation \(x_{ijk}\) is the value for sample \(i\), token position \(j\),
and feature channel \(k\). The representation of one token is therefore
\begin{equation}
    \mathbf{x}_{ij}
    =
    \begin{bmatrix}
        x_{ij1} & x_{ij2} & \cdots & x_{ijC}
    \end{bmatrix}^{\mathsf T}
    \in\mathbb{R}^{C}.
    \label{eq:norm-token-vector}
\end{equation}
The normalization methods most relevant to Transformers calculate their
statistics over the feature dimension of this token vector. The batch and
sequence dimensions are retained, but they do not determine the statistics
used to normalize a particular token.

The statistical quantities used in normalization can be introduced for a
generic vector
\(\mathbf{x}=[x_1,\ldots,x_C]^{\mathsf T}\in\mathbb{R}^{C}\).
Its arithmetic mean is
\begin{equation}
    \mu
    =
    \frac{1}{C}\sum_{k=1}^{C}x_k.
    \label{eq:norm-mean}
\end{equation}
The mean measures the central location of the feature values. Subtracting
this value from every component gives the centered vector
\begin{equation}
    \widetilde{x}_k=x_k-\mu.
    \label{eq:norm-centered-value}
\end{equation}
The population variance across the \(C\) features is
\begin{equation}
    \sigma^2
    =
    \frac{1}{C}\sum_{k=1}^{C}(x_k-\mu)^2,
    \label{eq:norm-variance}
\end{equation}
and the standard deviation is its nonnegative square root,
\begin{equation}
    \sigma=\sqrt{\sigma^2}.
    \label{eq:norm-standard-deviation}
\end{equation}
The factor \(1/C\), rather than \(1/(C-1)\), is used because the features
constitute the complete set over which the normalization statistics are
calculated; the expression is not an estimate of a population variance from
a separate statistical sample.

Subtracting the mean and dividing by the standard deviation produces the
standardized value
\begin{equation}
    \widehat{x}_k
    =
    \frac{x_k-\mu}{\sigma}.
    \label{eq:norm-standardized-value}
\end{equation}
When \(\sigma>0\), this transformation has zero mean and unit variance over
the feature dimension. In numerical implementations, the denominator is
modified by a small positive constant \(\varepsilon\):
\begin{equation}
    \widehat{x}_k
    =
    \frac{x_k-\mu}{\sqrt{\sigma^2+\varepsilon}}.
    \label{eq:norm-epsilon-standardization}
\end{equation}
The constant \(\varepsilon>0\) prevents division by zero when all feature
values are identical and reduces numerical sensitivity when the variance is
very small. With nonzero \(\varepsilon\), the resulting variance is not
exactly one; it is
\(\sigma^2/(\sigma^2+\varepsilon)\), which approaches one when
\(\varepsilon\) is negligible relative to \(\sigma^2\).

Most normalization operations can be expressed using the common affine
formulation
\begin{equation}
    \operatorname{Norm}(\mathbf{x})
    =
    \boldsymbol{\gamma}\odot
    \frac{\mathbf{x}-\boldsymbol{\mu}}
    {\sqrt{\boldsymbol{\sigma}^{2}+\varepsilon}}
    +
    \boldsymbol{\beta},
    \label{eq:norm-general-form}
\end{equation}
where the mean and variance are broadcast to the shape of the input,
\(\boldsymbol{\gamma}\in\mathbb{R}^{C}\) and
\(\boldsymbol{\beta}\in\mathbb{R}^{C}\) are learnable feature-wise scale
and shift parameters, and \(\odot\) denotes element-wise multiplication.
The scale and shift parameters allow the network to learn the output range
required by subsequent layers. Thus, normalization does not permanently
constrain the final output to have zero mean and unit variance; those
properties describe the standardized intermediate values before the affine
transformation. This common formulation and the dependence of different
normalization methods on the dimensions used to calculate their statistics
are described by \textcite{zhu2025dyt}.

The choice of aggregation dimensions distinguishes normalization methods.
For example, Batch Normalization calculates channel-wise statistics across
the batch and token dimensions. For channel \(k\), its mean and variance
are
\begin{align}
    \mu_k^{\mathrm{BN}}
    &=
    \frac{1}{BT}
    \sum_{i=1}^{B}\sum_{j=1}^{T}x_{ijk},
    \label{eq:norm-bn-mean}
    \\
    (\sigma_k^{\mathrm{BN}})^2
    &=
    \frac{1}{BT}
    \sum_{i=1}^{B}\sum_{j=1}^{T}
    (x_{ijk}-\mu_k^{\mathrm{BN}})^2.
    \label{eq:norm-bn-variance}
\end{align}
Other methods used in convolutional architectures, including Group
Normalization and Instance Normalization, use different groupings and
aggregation dimensions. They share the broad objective of transforming
activations using statistical quantities, but the axes over which those
quantities are calculated differ. For Transformer architectures, the two
principal statistical methods considered here are Layer Normalization and
Root Mean Square Normalization (RMSNorm)
\parencite{zhu2025dyt,corradini2025slmreview}.

\subsection{Layer Normalization: formulation and geometry}
\label{subsec:norm-layernorm}

Layer Normalization (LayerNorm or LN) calculates the mean and variance
independently for each token in each sample. For the token vector
\(\mathbf{x}_{ij}\), its mean is
\begin{equation}
    \mu_{ij}
    =
    \frac{1}{C}\sum_{k=1}^{C}x_{ijk},
    \label{eq:norm-ln-mean}
\end{equation}
and its variance is
\begin{equation}
    \sigma_{ij}^{2}
    =
    \frac{1}{C}\sum_{k=1}^{C}
    (x_{ijk}-\mu_{ij})^2.
    \label{eq:norm-ln-variance}
\end{equation}
The normalized output for feature \(k\), including the learnable affine
parameters, is
\begin{equation}
    \boxed{
    y_{ijk}^{\mathrm{LN}}
    =
    \gamma_k
    \frac{x_{ijk}-\mu_{ij}}
    {\sqrt{\sigma_{ij}^{2}+\varepsilon}}
    +\beta_k.
    }
    \label{eq:norm-layernorm}
\end{equation}
Here, \(\gamma_k\) and \(\beta_k\) are shared across samples and token
positions, but each feature channel has its own scale and shift parameters.
The statistics, in contrast, are calculated separately for every pair
\((i,j)\). The normalized representation of one token therefore does not
depend on the statistics of another token or another sample in the batch.

The distinction between the statistical transformation and the affine
parameters is essential. Before the affine transformation, the feature
values have zero mean and approximately unit variance. After multiplication
by \(\gamma_k\) and addition of \(\beta_k\), these properties need not hold.
The learnable parameters restore flexibility to the representation and allow
subsequent attention or feed-forward operations to receive feature values
with scales and offsets selected through training.

\paragraph{Geometric interpretation.}
The centering operation can be expressed as a projection in
\(\mathbb{R}^{C}\). Let \(\mathbf{1}\in\mathbb{R}^{C}\) be the all-ones
vector and let \(\mathbf{I}\) be the \(C\times C\) identity matrix. Define
the centering matrix
\begin{equation}
    \mathbf{P}
    =
    \mathbf{I}
    -
    \frac{1}{C}\mathbf{1}\mathbf{1}^{\mathsf T}.
    \label{eq:norm-centering-matrix}
\end{equation}
The mean of \(\mathbf{x}\) is
\(\mu=C^{-1}\mathbf{1}^{\mathsf T}\mathbf{x}\), so the centered vector is
\begin{equation}
    \mathbf{x}_{c}
    =
    \mathbf{P}\mathbf{x}.
    \label{eq:norm-projected-vector}
\end{equation}
The matrix \(\mathbf{P}\) is symmetric and idempotent:
\begin{equation}
    \mathbf{P}^{\mathsf T}=\mathbf{P},
    \qquad
    \mathbf{P}^{2}=\mathbf{P}.
    \label{eq:norm-projection-properties}
\end{equation}
It projects the vector onto the subspace orthogonal to \(\mathbf{1}\), whose
dimension is \(C-1\) for \(C>1\). The variance can consequently be written
as
\begin{equation}
    \sigma^2
    =
    \frac{1}{C}\mathbf{x}_{c}^{\mathsf T}\mathbf{x}_{c}
    =
    \frac{1}{C}\|\mathbf{P}\mathbf{x}\|_2^2.
    \label{eq:norm-variance-norm}
\end{equation}
Ignoring \(\varepsilon\), the standardized vector is
\begin{equation}
    \widehat{\mathbf{x}}
    =
    \frac{\sqrt{C}\,\mathbf{P}\mathbf{x}}
    {\|\mathbf{P}\mathbf{x}\|_2},
    \label{eq:norm-ln-geometric}
\end{equation}
provided that \(\mathbf{P}\mathbf{x}\neq\mathbf{0}\). It therefore lies on
a sphere of radius \(\sqrt{C}\) inside the zero-mean subspace. LayerNorm
first removes the component parallel to the all-ones vector and then fixes
the magnitude of the remaining centered vector.

This representation also reveals two invariance properties of the
statistical part of LayerNorm. For \(a>0\) and scalar \(b\), and when
\(\varepsilon\) is neglected,
\begin{equation}
    \operatorname{LN}_{\mathrm{stat}}
    (a\mathbf{x}+b\mathbf{1})
    =
    \operatorname{LN}_{\mathrm{stat}}(\mathbf{x}).
    \label{eq:norm-ln-invariance}
\end{equation}
Adding the same scalar to all features is removed by centering. Multiplying
all features by a positive scalar changes the centered vector and its
standard deviation by the same factor, so their ratio is unchanged. The
learnable affine transformation is not included in
\(\operatorname{LN}_{\mathrm{stat}}\) in
\cref{eq:norm-ln-invariance}.

\paragraph{Differential behavior.}
The Jacobian of a transformation describes how small changes in its input
affect its output. It is therefore central to understanding how gradients
are propagated through normalization during backpropagation. Define
\begin{equation}
    s
    =
    \sqrt{
        \frac{1}{C}\mathbf{x}_{c}^{\mathsf T}\mathbf{x}_{c}
        +\varepsilon
    },
    \qquad
    \widehat{\mathbf{x}}=\frac{\mathbf{x}_{c}}{s}.
    \label{eq:norm-ln-s-definition}
\end{equation}
Since \(\mathbf{x}_{c}=\mathbf{P}\mathbf{x}\), differentiating the
normalized vector with respect to the input gives
\begin{equation}
    \boxed{
    \frac{\partial\widehat{\mathbf{x}}}{\partial\mathbf{x}}
    =
    \frac{\mathbf{P}}{s}
    -
    \frac{\mathbf{x}_{c}\mathbf{x}_{c}^{\mathsf T}}
    {Cs^{3}}.
    }
    \label{eq:norm-ln-jacobian}
\end{equation}
The first term represents centering and scaling of input perturbations. The
second term accounts for the fact that the denominator also depends on the
input vector. A change in one feature can therefore alter the normalized
values of the other features through the shared mean and variance.

Let
\(\boldsymbol{\Gamma}
=\operatorname{diag}(\gamma_1,\ldots,\gamma_C)\).
The Jacobian of the complete affine LayerNorm operation is
\begin{equation}
    \frac{\partial\mathbf{y}^{\mathrm{LN}}}{\partial\mathbf{x}}
    =
    \boldsymbol{\Gamma}
    \left(
        \frac{\mathbf{P}}{s}
        -
        \frac{\mathbf{x}_{c}\mathbf{x}_{c}^{\mathsf T}}
        {Cs^{3}}
    \right).
    \label{eq:norm-ln-affine-jacobian}
\end{equation}
The shift vector \(\boldsymbol{\beta}\) does not appear because its
derivative with respect to the input is zero. The feature-wise scale
parameters modify the gradient magnitude of the corresponding output
features. These expressions are direct mathematical consequences of the
LayerNorm formulation in \cref{eq:norm-layernorm}; they are not additional
empirical results reported by the source paper.

\subsection{RMSNorm and normalization within residual Transformer blocks}
\label{subsec:norm-rmsnorm-residual}

Root Mean Square Normalization (RMSNorm) simplifies LayerNorm by removing
the mean-centering step. Instead of calculating squared deviations from the
mean, it calculates the mean of the squared feature values. For a token
representation \(\mathbf{x}_{ij}\), the mean square is
\begin{equation}
    m_{ij}^{(2)}
    =
    \frac{1}{C}\sum_{k=1}^{C}x_{ijk}^{2},
    \label{eq:norm-rms-mean-square}
\end{equation}
and the root mean square is
\begin{equation}
    r_{ij}
    =
    \sqrt{
        \frac{1}{C}\sum_{k=1}^{C}x_{ijk}^{2}
    }.
    \label{eq:norm-rms}
\end{equation}
The method can be expressed within the common affine formulation by setting
the mean to zero and using the mean square as the variance-like quantity:
\begin{equation}
    \boxed{
    y_{ijk}^{\mathrm{RMS}}
    =
    \gamma_k
    \frac{x_{ijk}}
    {\sqrt{
        \frac{1}{C}\sum_{\ell=1}^{C}x_{ij\ell}^{2}
        +\varepsilon
    }}
    +\beta_k.
    }
    \label{eq:norm-rmsnorm}
\end{equation}
The defining distinction is that RMSNorm does not subtract the token's
feature mean. The common affine expression in
\cref{eq:norm-general-form} is the parameterization used for the
normalization formulation in \textcite{zhu2025dyt}; implementations may
use a scale-only affine parameterization. The statistical distinction
between LayerNorm and RMSNorm is the use of the root mean square without
mean-centering, as also emphasized by
\textcite{corradini2025slmreview}.

The relation between variance and mean square follows by expanding the
definition of variance:
\begin{align}
    \frac{1}{C}\sum_{k=1}^{C}x_k^2
    &=
    \frac{1}{C}\sum_{k=1}^{C}
    \bigl[(x_k-\mu)+\mu\bigr]^2
    \nonumber\\
    &=
    \frac{1}{C}\sum_{k=1}^{C}(x_k-\mu)^2
    +\mu^2
    \nonumber\\
    &=
    \sigma^2+\mu^2.
    \label{eq:norm-mean-square-variance}
\end{align}
Thus, mean square and variance are equal only when the feature mean is zero.
RMSNorm controls the root mean square magnitude of the vector, but it does
not explicitly force the mean of its features to zero.

\paragraph{Geometric interpretation and Jacobian.}
Ignoring \(\varepsilon\), the root mean square is
\(r=\|\mathbf{x}\|_2/\sqrt{C}\), and the normalized vector is
\begin{equation}
    \widehat{\mathbf{x}}^{\mathrm{RMS}}
    =
    \frac{\sqrt{C}}{\|\mathbf{x}\|_2}\mathbf{x}.
    \label{eq:norm-rms-geometric}
\end{equation}
Its Euclidean norm is \(\sqrt{C}\), but it is not restricted to the
zero-mean subspace. RMSNorm rescales the original vector without first
projecting out its component in the direction of \(\mathbf{1}\). It is
invariant to positive scalar multiplication when \(\varepsilon\) is
neglected, but it is not invariant to adding a common scalar to all
features.

For
\begin{equation}
    r
    =
    \sqrt{
        \frac{1}{C}\mathbf{x}^{\mathsf T}\mathbf{x}
        +\varepsilon
    },
    \qquad
    \widehat{\mathbf{x}}^{\mathrm{RMS}}=\frac{\mathbf{x}}{r},
    \label{eq:norm-rms-r-definition}
\end{equation}
the Jacobian is
\begin{equation}
    \boxed{
    \frac{\partial\widehat{\mathbf{x}}^{\mathrm{RMS}}}
    {\partial\mathbf{x}}
    =
    \frac{\mathbf{I}}{r}
    -
    \frac{\mathbf{x}\mathbf{x}^{\mathsf T}}{Cr^{3}}.
    }
    \label{eq:norm-rms-jacobian}
\end{equation}
As in LayerNorm, the denominator couples the feature values of a token:
changing one feature changes the root mean square used to normalize all
features. The difference is that RMSNorm does not contain the centering
projection \(\mathbf{P}\).

\paragraph{Residual connections.}
Normalization is used together with residual connections in Transformer
blocks. Let \(\mathbf{H}\) be the input to a sublayer and let \(F\) denote
the transformation performed by that sublayer, such as an attention or
feed-forward operation. A residual addition forms
\begin{equation}
    \mathbf{Z}
    =
    \mathbf{H}+F(\mathbf{H}).
    \label{eq:norm-residual-addition}
\end{equation}
The derivative of this addition with respect to its input is
\begin{equation}
    \frac{\partial\mathbf{Z}}{\partial\mathbf{H}}
    =
    \mathbf{I}+\mathbf{J}_{F},
    \label{eq:norm-residual-jacobian}
\end{equation}
where \(\mathbf{J}_{F}\) is the Jacobian of \(F\). The identity term
provides a direct contribution to the gradient independently of the
derivative of the sublayer transformation.

In the post-residual arrangement described in the systematic review,
normalization is applied after the
sublayer output has been added to its input:
\begin{equation}
    \mathbf{H}_{\mathrm{next}}
    =
    \operatorname{LN}
    \bigl(\mathbf{H}+F(\mathbf{H})\bigr).
    \label{eq:norm-post-residual}
\end{equation}
An RMSNorm operation can occupy the corresponding normalization position
in architectures that use RMSNorm. Residual addition and normalization
therefore have complementary roles in this described block arrangement:
the residual path provides a direct route for information and gradients,
while normalization controls the numerical scale of the representation
passed to subsequent computation. The equation specifies the arrangement
discussed by the review and should not be read as a claim that every
Transformer implementation uses an identical normalization placement.

\begin{figure}[htbp]
    \centering
    % Figure source fragment: comparison of LayerNorm, RMSNorm, and DyT.
% This file is inserted inside a figure environment using \input.
\begin{center}
\small
\setlength{\fboxsep}{6pt}
\begin{tabular}{ccc}
    \textbf{LayerNorm} &
    \textbf{RMSNorm} &
    \textbf{Dynamic Tanh} \\[0.4em]
    \fbox{\parbox{0.28\textwidth}{\centering
        \textbf{Input token}\\
        \(\mathbf{x}\in\mathbb{R}^{C}\)\\[0.25em]
        \(\downarrow\)\\[0.2em]
        \textbf{Token statistics}\\
        \(\mu=\frac{1}{C}\sum_k x_k\)\\
        \(\sigma^2=\frac{1}{C}\sum_k(x_k-\mu)^2\)\\[0.2em]
        \(\downarrow\)\\[0.2em]
        \textbf{Center and scale}\\
        \(\widehat{\mathbf{x}}=
        \frac{\mathbf{x}-\mu}{\sqrt{\sigma^2+\varepsilon}}\)\\[0.2em]
        \(\downarrow\)\\[0.2em]
        \textbf{Affine output}\\
        \(\boldsymbol{\gamma}\odot\widehat{\mathbf{x}}+
        \boldsymbol{\beta}\)}} &
    \fbox{\parbox{0.28\textwidth}{\centering
        \textbf{Input token}\\
        \(\mathbf{x}\in\mathbb{R}^{C}\)\\[0.25em]
        \(\downarrow\)\\[0.2em]
        \textbf{Token RMS}\\
        \(r=\sqrt{\frac{1}{C}\sum_k x_k^2}\)\\[0.2em]
        \(\downarrow\)\\[0.2em]
        \textbf{Scale without centering}\\
        \(\widehat{\mathbf{x}}=\frac{\mathbf{x}}{\sqrt{r^2+\varepsilon}}\)\\[0.2em]
        \(\downarrow\)\\[0.2em]
        \textbf{Scaled output}\\
        \(\boldsymbol{\gamma}\odot\widehat{\mathbf{x}}\)}} &
    \fbox{\parbox{0.28\textwidth}{\centering
        \textbf{Input token}\\
        \(\mathbf{x}\in\mathbb{R}^{C}\)\\[0.25em]
        \(\downarrow\)\\[0.2em]
        \textbf{Learnable scaling}\\
        \(\alpha\mathbf{x}\)\\[0.2em]
        \(\downarrow\)\\[0.2em]
        \textbf{Element-wise tanh}\\
        \(u_k=\tanh(\alpha x_k)\)\\[0.2em]
        \(\downarrow\)\\[0.2em]
        \textbf{Affine output}\\
        \(\boldsymbol{\gamma}\odot\mathbf{u}+
        \boldsymbol{\beta}\)}}
\end{tabular}
\end{center}

    \caption{Schematic comparison of the operations performed by LayerNorm,
    RMSNorm, and Dynamic Tanh. LayerNorm calculates token-wise mean and
    variance and centers the feature vector; RMSNorm calculates the
    token-wise root mean square without centering; DyT applies a learnable
    element-wise scaling followed by a hyperbolic tangent. The comparison
    follows the formulations in \textcite{zhu2025dyt} and
    \textcite{corradini2025slmreview}.}
    \label{fig:norm-method-comparison}
\end{figure}
% Alt text for accessibility register: Three parallel columns compare
% LayerNorm, RMSNorm, and Dynamic Tanh. The LayerNorm column calculates a
% token mean and variance, centers and scales the vector, and applies
% channel-wise affine parameters. The RMSNorm column calculates the root
% mean square and rescales without centering. The Dynamic Tanh column
% multiplies each activation by a learnable scalar, applies tanh, and then
% applies channel-wise scale and shift parameters.

\subsection{Observed input-output behavior of trained normalization layers}
\label{subsec:norm-observed-behavior}

The statistical definitions explain how LayerNorm and RMSNorm are computed,
but they do not by themselves describe the aggregate input-output behavior
observed across a trained Transformer. To investigate this behavior,
\textcite{zhu2025dyt} examine normalization layers in three trained
architectures: a Vision Transformer (ViT-B) trained on ImageNet-1K, a
wav2vec 2.0 Large Transformer trained on LibriSpeech, and a Diffusion
Transformer (DiT-XL) trained on ImageNet-1K. In the examined models,
LayerNorm is applied within the Transformer blocks and before the final
linear projection.

The authors perform forward passes on mini-batches and record the tensor
elements immediately before and after selected LayerNorm operations,
before the learnable affine transformation. Because LayerNorm preserves
the tensor dimensions, each input activation can be paired with its
corresponding normalized output. Plotting these pairs reveals that the
relationships in earlier normalization layers are predominantly linear,
whereas deeper layers often exhibit full or partial S-shaped curves.

An important distinction is required when interpreting these plots. For
one token, the mean and standard deviation are fixed across its feature
channels during the normalization operation. Its input-output relation is
therefore
\begin{equation}
    y_k
    =
    \frac{x_k-\mu}{\sqrt{\sigma^2+\varepsilon}},
    \label{eq:norm-token-affine-map}
\end{equation}
which is affine in \(x_k\) for that token. Different tokens, however, have
different means and standard deviations. Their token-specific affine
relations consequently have different slopes and intercepts. When
activations from many tokens are plotted together, the collection of these
different lines can form an aggregate nonlinear pattern.

The token-grouped visualization reported by \textcite{zhu2025dyt} shows
approximately straight relationships within individual tokens, with
different slopes associated with different token variances. The channel-
grouped visualization shows that channels can span substantially different
input ranges. A small number of channels contain extreme values, and these
values are strongly reduced by the normalization operation. The aggregate
input-output relationship therefore resembles a scaled hyperbolic tangent,
even though the operation applied to each individual token is based on
centering and division by its standard deviation.

The authors report that approximately \(99\%\) of the activation points in
their analysis fall within the central, approximately linear part of the
observed S-shaped relationship. Nevertheless, some activations are
extreme; in the examined ViT, examples include values greater than \(50\)
or smaller than \(-50\). The normalization operation reduces the relative
magnitude of these outlying values and brings them closer to the range
occupied by the majority of activations. This disproportionate treatment
of extreme values is not reproduced by a single affine transformation with
a constant slope. The observation motivates the hypothesis that nonlinear
suppression of extreme activations is an important functional effect of
normalization in the examined trained networks
\parencite{zhu2025dyt}.

The hyperbolic tangent function provides a simple mathematical model of
the observed S-shaped relation. It is defined by
\begin{equation}
    \tanh(x)
    =
    \frac{e^x-e^{-x}}{e^x+e^{-x}},
    \label{eq:norm-tanh-definition}
\end{equation}
and maps real inputs into the bounded interval \((-1,1)\). Its limiting
behavior is
\begin{equation}
    \lim_{x\to+\infty}\tanh(x)=1,
    \qquad
    \lim_{x\to-\infty}\tanh(x)=-1,
    \qquad
    \tanh(0)=0.
    \label{eq:norm-tanh-limits}
\end{equation}
Near the origin, its Taylor expansion is
\begin{equation}
    \tanh(x)
    =
    x-\frac{x^3}{3}
    +\frac{2x^5}{15}
    -\frac{17x^7}{315}
    +\cdots,
    \label{eq:norm-tanh-taylor}
\end{equation}
so the function is approximately linear when \(|x|\) is small. Its
derivative is
\begin{equation}
    \frac{d}{dx}\tanh(x)
    =
    1-\tanh^2(x)
    =
    \operatorname{sech}^2(x).
    \label{eq:norm-tanh-derivative}
\end{equation}
The derivative is largest at the origin and approaches zero as the
magnitude of the input increases. Thus, the central region transmits
small changes approximately linearly, while the saturated regions reduce
the sensitivity of the output to further increases in the input magnitude.

A scaled hyperbolic tangent,
\begin{equation}
    f_{\alpha}(x)=\tanh(\alpha x),
    \label{eq:norm-scaled-tanh}
\end{equation}
changes the width of the transition region. For \(\alpha>0\), increasing
\(\alpha\) causes saturation to occur over a narrower input range, whereas
decreasing \(\alpha\) expands the approximately linear region. This
scaling property is the mathematical basis for the learnable parameter
introduced by Dynamic Tanh.

\begin{figure}[htbp]
    \centering
    % Figure source fragment: schematic hyperbolic-tangent saturation curve.
% This file is inserted inside a figure environment using \input.
% The curve is schematic and is not a plot of experimental measurements.
\begin{center}
\setlength{\unitlength}{0.68cm}
\begin{picture}(8.4,4.4)
    % Coordinate axes.
    \put(0.45,2.15){\vector(1,0){7.45}}
    \put(4.10,0.30){\vector(0,1){3.80}}

    % Approximate S-shaped curve using quadratic Bezier segments.
    \qbezier(0.75,0.55)(1.55,0.55)(2.35,0.92)
    \qbezier(2.35,0.92)(3.10,1.38)(3.55,1.83)
    \qbezier(3.55,1.83)(3.83,2.08)(4.10,2.15)
    \qbezier(4.10,2.15)(4.37,2.22)(4.65,2.47)
    \qbezier(4.65,2.47)(5.10,2.92)(5.85,3.38)
    \qbezier(5.85,3.38)(6.65,3.75)(7.45,3.75)

    % Saturation guides.
    \put(0.75,3.75){\line(1,0){6.70}}
    \put(0.75,0.55){\line(1,0){6.70}}

    % Axis labels and function annotation.
    \put(7.95,1.93){\(x\)}
    \put(4.23,4.02){\(y\)}
    \put(7.10,3.86){\(1\)}
    \put(7.10,0.20){\(-1\)}
    \put(4.22,1.88){\(0\)}
    \put(1.05,2.38){\(\tanh(\alpha x)\)}
\end{picture}
\end{center}

    \caption{Schematic S-shaped input-output relation of the hyperbolic
    tangent. The central region is approximately linear, while the output
    approaches bounded positive and negative limits for large input
    magnitudes. The role of the scale parameter is discussed by
    \textcite{zhu2025dyt}.}
    \label{fig:norm-tanh-curve}
\end{figure}
% Alt text for accessibility register: A two-dimensional plot of tanh(alpha x)
% with an increasing S-shaped curve, horizontal saturation levels near -1
% and +1, and a central region around the origin that is approximately linear.

\subsection{Dynamic Tanh: formulation, derivatives, and implementation}
\label{subsec:norm-dyt-formulation}

The S-shaped aggregate input-output relationships observed in trained
normalization layers motivate an alternative to statistical normalization.
Rather than calculating the mean and variance of each token representation,
Dynamic Tanh (DyT) applies a learnable element-wise nonlinear transformation.
For an input tensor \(\mathbf{X}\), the operation introduced by
\textcite{zhu2025dyt} is
\begin{equation}
    \boxed{
    \operatorname{DyT}(\mathbf{X})
    =
    \boldsymbol{\gamma}\odot
    \tanh(\alpha\mathbf{X})
    +
    \boldsymbol{\beta}.
    }
    \label{eq:norm-dyt}
\end{equation}
The scalar \(\alpha\) is learnable and shared across the elements of a DyT
layer. The vectors
\(\boldsymbol{\gamma},\boldsymbol{\beta}\in\mathbb{R}^{C}\) are
learnable channel-wise scale and shift parameters. The hyperbolic tangent
is applied element-wise, and the channel-wise parameters are broadcast
across the batch and token dimensions. For one activation, the operation
is
\begin{equation}
    y_{ijk}
    =
    \gamma_k\tanh(\alpha x_{ijk})+\beta_k.
    \label{eq:norm-dyt-component}
\end{equation}

The scalar \(\alpha\) controls the scale of the input before the nonlinear
function. The parameter \(\gamma_k\) controls the output scale of channel
\(k\), and \(\beta_k\) controls its output displacement. If \(\alpha>0\),
the output of channel \(k\) approaches
\(\beta_k+\gamma_k\) for large positive inputs and
\(\beta_k-\gamma_k\) for large negative inputs. More generally, its
bounded range is determined by \(\beta_k\) and \(|\gamma_k|\). Unlike
LayerNorm and RMSNorm, DyT does not calculate a mean, variance, or root
mean square from the current input tensor.

\paragraph{Local behavior and saturation.}
For one feature, define
\begin{equation}
    y_k
    =
    \gamma_k\tanh(\alpha x_k)+\beta_k.
    \label{eq:norm-dyt-scalar}
\end{equation}
Its derivative with respect to the input is
\begin{equation}
    \boxed{
    \frac{\partial y_k}{\partial x_k}
    =
    \gamma_k\alpha
    \left[1-\tanh^2(\alpha x_k)\right].
    }
    \label{eq:norm-dyt-input-derivative}
\end{equation}
At the origin, the derivative is \(\gamma_k\alpha\). For small
\(|\alpha x_k|\), the Taylor expansion gives
\begin{equation}
    y_k
    \approx
    \gamma_k\alpha x_k+\beta_k.
    \label{eq:norm-dyt-local-linear}
\end{equation}
For large \(|\alpha x_k|\), the derivative approaches zero and the output
saturates. This supplies an element-wise mechanism for reducing the
sensitivity of the output to extreme input values.

The derivative with respect to the shared scale parameter is
\begin{equation}
    \frac{\partial y_k}{\partial\alpha}
    =
    \gamma_k x_k
    \left[1-\tanh^2(\alpha x_k)\right].
    \label{eq:norm-dyt-alpha-derivative}
\end{equation}
For a scalar loss \(\mathcal{L}\), the parameter gradients accumulate
contributions from the batch and sequence dimensions:
\begin{align}
    \frac{\partial\mathcal{L}}{\partial\beta_k}
    &=
    \sum_{i,j}
    \frac{\partial\mathcal{L}}{\partial y_{ijk}},
    \label{eq:norm-dyt-beta-gradient}
    \\
    \frac{\partial\mathcal{L}}{\partial\gamma_k}
    &=
    \sum_{i,j}
    \frac{\partial\mathcal{L}}{\partial y_{ijk}}
    \tanh(\alpha x_{ijk}),
    \label{eq:norm-dyt-gamma-gradient}
    \\
    \frac{\partial\mathcal{L}}{\partial\alpha}
    &=
    \sum_{i,j,k}
    \frac{\partial\mathcal{L}}{\partial y_{ijk}}
    \gamma_k x_{ijk}
    \left[1-\tanh^2(\alpha x_{ijk})\right].
    \label{eq:norm-dyt-alpha-gradient}
\end{align}
These expressions show how the scalar parameter is learned from the
activation elements throughout the layer. Contributions from highly
saturated activations are reduced by the factor
\(1-\tanh^2(\alpha x_{ijk})\).

For a token vector, define
\(\boldsymbol{\Gamma}
=\operatorname{diag}(\gamma_1,\ldots,\gamma_C)\).
The Jacobian of DyT is diagonal:
\begin{equation}
    \boxed{
    \mathbf{J}_{\mathrm{DyT}}
    =
    \boldsymbol{\Gamma}
    \operatorname{diag}
    \left(
        \alpha\left[
            1-\tanh^2(\alpha\mathbf{x})
        \right]
    \right).
    }
    \label{eq:norm-dyt-jacobian}
\end{equation}
The expression inside the second diagonal matrix is understood
element-wise. The diagonal structure follows from the fact that each
output activation depends directly only on its corresponding input
activation. This differs from LayerNorm and RMSNorm, whose shared
token-level denominators couple the feature values in the Jacobian.
The derivative comparison is a mathematical consequence of the
respective forward equations.

\paragraph{Implementation and initialization.}
For an input tensor of shape \([B,T,C]\), the implementation in
\textcite{zhu2025dyt} initializes the scalar parameter to a chosen value
\(\alpha_0\), initializes the channel-wise scale vector to ones, and
initializes the channel-wise shift vector to zeros:
\begin{equation}
    \alpha=\alpha_0,
    \qquad
    \boldsymbol{\gamma}=\mathbf{1}_C,
    \qquad
    \boldsymbol{\beta}=\mathbf{0}_C.
    \label{eq:norm-dyt-initialization}
\end{equation}
The forward pass first computes
\begin{equation}
    \mathbf{U}=\tanh(\alpha\mathbf{X}),
    \label{eq:norm-dyt-intermediate}
\end{equation}
and then applies
\begin{equation}
    \mathbf{Y}
    =
    \boldsymbol{\gamma}\odot\mathbf{U}
    +\boldsymbol{\beta}.
    \label{eq:norm-dyt-output}
\end{equation}
The operation replaces one normalization layer with one DyT layer. It
can be inserted at normalization locations in attention blocks,
feed-forward blocks, and before the final linear projection. The original
activation functions in the architecture, such as GELU or ReLU, are not
replaced by DyT; the method replaces the normalization operation while
leaving those activation functions unchanged
\parencite{zhu2025dyt}.

The learnable scalar \(\alpha\) is not calculated from the current token's
statistics. The analysis in \textcite{zhu2025dyt} nevertheless finds that
its value tracks the inverse standard deviation of the input activations
during training and is correlated with that quantity in trained models.
This indicates that \(\alpha\) can partially learn an activation-scale
adjustment. However, it is shared across the input activations of a layer,
rather than being calculated separately for every token. The scalar alone
therefore cannot reproduce token-specific statistical normalization or
the nonlinear suppression of extreme values; the latter is supplied by
the hyperbolic tangent function.

\subsection{Computational analysis and a worked numerical example}
\label{subsec:norm-computation-example}

The three methods differ in the operations required for each token
representation. LayerNorm calculates the feature mean, subtracts it,
calculates the mean squared centered deviation, adds the stability
constant, takes a square root, and scales the centered features. RMSNorm
calculates the mean of squared features, adds the stability constant,
takes a square root, and scales the original features. DyT performs a
scalar multiplication, an element-wise hyperbolic tangent, and a
channel-wise affine transformation.

For a tensor with \(B T\) token representations and \(C\) features per
token, each method requires work proportional to the number of activation
elements. Their leading arithmetic complexity is therefore
\begin{equation}
    \mathcal{O}(BTC).
    \label{eq:norm-complexity}
\end{equation}
This asymptotic expression does not imply equal execution time. LayerNorm
and RMSNorm require reductions over the feature dimension, whereas DyT
uses element-wise operations and a nonlinear function. Actual latency
depends on hardware, memory traffic, kernel fusion, numerical precision,
and compiler optimization. In particular, \textcite{zhu2025dyt} report
that DyT does not provide a speedup over properly compiled and optimized
normalization layers, despite its simpler conceptual formulation.

A numerical example makes the transformations explicit. Consider the
four-dimensional token vector
\begin{equation}
    \mathbf{x}
    =
    \begin{bmatrix}
        1 & 2 & 3 & 4
    \end{bmatrix}^{\mathsf T},
    \qquad C=4.
    \label{eq:norm-example-input}
\end{equation}
For this example, set \(\varepsilon=0\) to keep the arithmetic transparent,
and set the affine scale to one and the affine shift to zero.

The LayerNorm mean is
\begin{equation}
    \mu
    =
    \frac{1+2+3+4}{4}
    =
    2.5.
    \label{eq:norm-example-mean}
\end{equation}
The variance is
\begin{align}
    \sigma^2
    &=
    \frac{(1-2.5)^2+(2-2.5)^2
    +(3-2.5)^2+(4-2.5)^2}{4}
    \nonumber\\
    &=
    \frac{2.25+0.25+0.25+2.25}{4}
    =
    1.25,
    \label{eq:norm-example-variance}
\end{align}
and the standard deviation is
\begin{equation}
    \sigma=\sqrt{1.25}\approx1.1180.
    \label{eq:norm-example-std}
\end{equation}
The resulting LayerNorm vector is approximately
\begin{equation}
    \widehat{\mathbf{x}}^{\mathrm{LN}}
    =
    \begin{bmatrix}
        -1.3416 & -0.4472 & 0.4472 & 1.3416
    \end{bmatrix}^{\mathsf T}.
    \label{eq:norm-example-ln-output}
\end{equation}

For RMSNorm, the mean square is
\begin{equation}
    m^{(2)}
    =
    \frac{1^2+2^2+3^2+4^2}{4}
    =
    \frac{30}{4}
    =
    7.5,
    \label{eq:norm-example-mean-square}
\end{equation}
so that
\begin{equation}
    r=\sqrt{7.5}\approx2.7386.
    \label{eq:norm-example-rms}
\end{equation}
The RMS-normalized vector is consequently
\begin{equation}
    \widehat{\mathbf{x}}^{\mathrm{RMS}}
    \approx
    \begin{bmatrix}
        0.3651 & 0.7303 & 1.0954 & 1.4606
    \end{bmatrix}^{\mathsf T}.
    \label{eq:norm-example-rms-output}
\end{equation}
Unlike LayerNorm, RMSNorm does not subtract the mean, so the signs and
relative direction of the original vector are retained.

For DyT, choose \(\alpha=0.5\), \(\boldsymbol{\gamma}=\mathbf{1}\),
and \(\boldsymbol{\beta}=\mathbf{0}\). Its output is
\begin{equation}
    \operatorname{DyT}(\mathbf{x})
    =
    \tanh(0.5\mathbf{x})
    \approx
    \begin{bmatrix}
        0.4621 & 0.7616 & 0.9051 & 0.9640
    \end{bmatrix}^{\mathsf T}.
    \label{eq:norm-example-dyt-output}
\end{equation}
The fourth component is close to the positive saturation limit, whereas
the first component remains in the more nearly linear region. This example
illustrates the different mathematical operations; it is not an
experimental result reported by either source paper.

\subsection{Language-model evaluation and initialization}
\label{subsec:norm-language-model-evaluation}

The replacement of a normalization layer is useful only if it preserves
the behavior of the complete model during training and evaluation.
\textcite{zhu2025dyt} investigate this question by replacing RMSNorm with
DyT in LLaMA models of different sizes. The reported configurations are
LLaMA 7B, 13B, 34B, and 70B. The models are pretrained on The Pile using
200 billion tokens and follow the LLaMA training recipe. The DyT
experiments also introduce a learnable scalar parameter after the initial
embedding layer and tune the initial value of \(\alpha\).

The authors report the final training loss and average performance over
15 zero-shot tasks from the lm-eval framework. The results are summarized
in \cref{tab:norm-llama-results}. The zero-shot scores are identical at
the reported precision for the RMSNorm and DyT versions of all four
models. The final losses differ by \(0.01\) for the 7B and 13B models and
are identical at the reported precision for the 34B and 70B models.
The training curves are also closely aligned across the four model sizes
\parencite{zhu2025dyt}.

\begin{table}[htbp]
    \centering
    \caption{LLaMA pretraining loss and average performance on 15 zero-shot
    lm-eval tasks when using RMSNorm or DyT. Values are reproduced from
    Table 4 of \textcite{zhu2025dyt}.}
    \label{tab:norm-llama-results}
    \begin{tabular}{lcccc}
        \toprule
        Model
        & RMSNorm score
        & DyT score
        & RMSNorm loss
        & DyT loss \\
        \midrule
        LLaMA 7B  & 0.513 & 0.513 & 1.59 & 1.60 \\
        LLaMA 13B & 0.529 & 0.529 & 1.53 & 1.54 \\
        LLaMA 34B & 0.536 & 0.536 & 1.50 & 1.50 \\
        LLaMA 70B & 0.549 & 0.549 & 1.45 & 1.45 \\
        \bottomrule
    \end{tabular}
\end{table}

\paragraph{The role of the initial value of \(\alpha\).}
The authors find that an initial value of \(\alpha_0=0.5\) is generally
sufficient for the non-language-model tasks examined in their study.
The situation differs for LLM training: tuning \(\alpha_0\) can
substantially improve the training loss. Their LLaMA experiments indicate
that larger models require smaller initial values and that DyT layers in
attention blocks benefit from larger initial values than those in
feed-forward blocks or before the final output projection.

The model-specific values reported in Table 10 of
\textcite{zhu2025dyt} are shown in
\cref{tab:norm-llama-alpha-init}. The distinction between the attention
and other locations becomes particularly important as model width
increases.

\begin{table}[htbp]
    \centering
    \caption{Reported initial values of the DyT scaling parameter for
    attention blocks and other DyT locations in LLaMA models. Values are
    reproduced from Table 10 of \textcite{zhu2025dyt}.}
    \label{tab:norm-llama-alpha-init}
    \begin{tabular}{lccc}
        \toprule
        Model & Width & Depth & \(\alpha_0\) (attention / other) \\
        \midrule
        LLaMA 7B  & 4096 & 32 & \(0.8/0.20\) \\
        LLaMA 13B & 5120 & 40 & \(0.6/0.15\) \\
        LLaMA 34B & 8196 & 48 & \(0.2/0.05\) \\
        LLaMA 70B & 8196 & 80 & \(0.2/0.05\) \\
        \bottomrule
    \end{tabular}
\end{table}

To separate the influence of model width and depth, the authors also
report optimal initializations across combinations of widths and depths.
The values are reproduced in
\cref{tab:norm-alpha-width-depth}. The selected initialization remains
constant across the tested depths at a given width, while it decreases
as width increases. The authors therefore identify width as the more
important factor in these experiments and report comparatively little
influence from depth.

\begin{table}[htbp]
    \centering
    \caption{Reported optimal \(\alpha_0\) values
    (attention / other locations) across model widths and depths.
    Values are reproduced from Table 11 of
    \textcite{zhu2025dyt}.}
    \label{tab:norm-alpha-width-depth}
    \begin{tabular}{lcccc}
        \toprule
        Width & Depth 8 & Depth 16 & Depth 32 & Depth 64 \\
        \midrule
        1024 & \(1.0/1.0\) & \(1.0/1.0\) & \(1.0/1.0\) & \(1.0/1.0\) \\
        2048 & \(1.0/0.5\) & \(1.0/0.5\) & \(1.0/0.5\) & \(1.0/0.5\) \\
        4096 & \(0.8/0.2\) & \(0.8/0.2\) & \(0.8/0.2\) & \(0.8/0.2\) \\
        8192 & \(0.2/0.05\) & \(0.2/0.05\) & \(0.2/0.05\) & \(0.2/0.05\) \\
        \bottomrule
    \end{tabular}
\end{table}

The dependence on \(\alpha_0\) is also related to training stability.
In the supervised ViT experiments, larger models are more prone to
instability, and reducing either the learning rate or the initial
\(\alpha_0\) improves stability. The authors report that a default
\(\alpha_0=0.5\) gives stability comparable to LayerNorm for the
non-LLM settings they examine. For LLMs, however, their results motivate
separate tuning of the initial values at different architectural
locations rather than relying on one universal value
\parencite{zhu2025dyt}.

\paragraph{Ablation of the squashing function and the learnable scale.}
The role of the hyperbolic tangent is tested by replacing it with
hardtanh, sigmoid, or the identity function. In the ViT-S and ViT-B
experiments, using the identity function leads to training divergence,
whereas the squashing functions permit stable training. Among the
tested squashing functions, tanh achieves the highest reported accuracy.
The results are reproduced in
\cref{tab:norm-squashing-ablation}.

\begin{table}[htbp]
    \centering
    \caption{ImageNet-1K classification accuracy with alternative DyT
    squashing functions. The identity-function experiments diverged;
    the preceding accuracy is the highest reached before divergence.
    Values are reproduced from Table 7 of
    \textcite{zhu2025dyt}.}
    \label{tab:norm-squashing-ablation}
    \begin{tabular}{lcccc}
        \toprule
        Model & Identity & Tanh & Hardtanh & Sigmoid \\
        \midrule
        ViT-S & \(58.5\%\rightarrow\text{failed}\)
              & \(80.3\%\) & \(79.9\%\) & \(79.6\%\) \\
        ViT-B & \(61.0\%\rightarrow\text{failed}\)
              & \(82.5\%\) & \(82.2\%\) & \(81.6\%\) \\
        \bottomrule
    \end{tabular}
\end{table}

The learnable scale is examined separately by comparing the same
squashing functions with and without \(\alpha\). For ViT-B, removing
\(\alpha\) reduces the reported accuracy for all three nonlinear
functions. With tanh, the accuracy decreases from \(82.5\%\) to
\(81.1\%\); with hardtanh, it decreases from \(82.2\%\) to \(80.7\%\);
and with sigmoid, it decreases from \(81.6\%\) to \(80.7\%\).
These results indicate that the squashing function and its learnable
input scale make complementary contributions in the tested setting
\parencite{zhu2025dyt}.

\subsection{Normalization choices in Small Language Models}
\label{subsec:norm-small-language-models}

The choice between LayerNorm and RMSNorm is also visible in the
architectural trends of Small Language Models. These models are designed
to balance language-model capability against constraints on parameter
count, memory, inference time, and deployment hardware. In this setting,
the computational structure of repeatedly applied Transformer
components becomes an important design consideration.

The systematic review by \textcite{corradini2025slmreview} analyzes
70 English-language studies published between January 2023 and January
2025, focusing on recent SLM developments and primarily on models with
up to seven billion parameters. Its architectural analysis identifies
RMSNorm as a frequently adopted component. The review reports RMSNorm
in approximately \(66\%\) of the analyzed model designs and describes
its prevalence as a preference for a computationally lighter alternative
to standard LayerNorm.

This prevalence is consistent with the mathematical distinction already
established. RMSNorm omits the mean-centering operation and uses the
root mean square of the token features to determine the scale. Repeating
this simpler statistical calculation across many tokens and layers can
be relevant to the computational design of compact language models.
However, the reported \(66\%\) is a frequency within the review's
selected research corpus; it is not a universal estimate for every
existing SLM or language model architecture. Nor does the frequency
alone establish a controlled performance advantage over LayerNorm
\parencite{corradini2025slmreview}.

The three methods therefore represent different choices about how
activation behavior is controlled. LayerNorm uses token-specific
centering and variance scaling. RMSNorm retains token-specific
statistical scaling but removes mean-centering. DyT does not calculate
token statistics and instead uses a learnable element-wise nonlinear
transformation. Their common architectural purpose does not make their
mathematical operations interchangeable in every setting.

\begin{table}[htbp]
    \centering
    \small
    \setlength{\tabcolsep}{3pt}
    \caption{Conceptual comparison of the three methods based on the
    formulations discussed by \textcite{zhu2025dyt} and
    \textcite{corradini2025slmreview}.}
    \label{tab:norm-method-comparison}
    \begin{tabularx}{\textwidth}{p{0.22\textwidth}XXX}
        \toprule
        Property & LayerNorm & RMSNorm & Dynamic Tanh \\
        \midrule
        Mean calculation
        & Yes
        & No
        & No \\
        Variance or RMS calculation
        & Variance
        & Root mean square
        & No input statistic \\
        Mean-centering
        & Yes
        & No
        & No \\
        Input scaling
        & Token-specific standard deviation
        & Token-specific root mean square
        & Learnable layer-wise scalar \\
        Nonlinear saturation
        & Indirect aggregate behavior
        & No explicit squashing function
        & Explicit hyperbolic tangent \\
        Feature coupling in the forward Jacobian
        & Yes
        & Yes
        & No; diagonal Jacobian \\
        \bottomrule
    \end{tabularx}
\end{table}

\subsection{Limitations and connection to the attention mechanism}
\label{subsec:norm-limitations-transition}

The results of \textcite{zhu2025dyt} challenge the assumption that
statistical normalization layers are indispensable for training modern
Transformer architectures. They do not, however, establish that DyT is
a universal replacement for every normalization method or architecture.
The experiments primarily investigate networks that use LayerNorm or
RMSNorm. The authors report that preliminary experiments do not support
a direct replacement of Batch Normalization by DyT in conventional
ResNet architectures. The applicability of DyT to other normalization
methods therefore remains unresolved.

The computational implications also require qualification. Although
DyT removes the statistical reductions used by LayerNorm and RMSNorm,
the experiments report no speedup over optimized normalization layers.
Actual latency depends on hardware and implementation; a simpler
expression does not guarantee lower deployment latency
\parencite{zhu2025dyt}.

The evidence from the two provided studies supports a more precise
understanding of normalization as an architectural design choice.
The systematic review identifies RMSNorm as a prevalent option in
recent SLM research, while the DyT study demonstrates that a learned,
bounded element-wise transformation can achieve comparable performance
to RMSNorm in the evaluated LLaMA configurations. These results do not
show that all normalization mechanisms have identical effects. Rather,
they show that the control of activation scale and the suppression of
extreme values can be implemented through different mathematical
operations.

Normalization is therefore best understood not merely as a procedure
for forcing activations to have a prescribed statistical distribution,
but as a component that affects the numerical behavior and
differential structure of the transformations composing a deep
network. Its interaction with residual pathways, attention sublayers,
and feed-forward transformations determines how intermediate
representations are passed through the architecture.

With the role of normalization established, the next architectural
component to examine is attention. Whereas normalization controls
aspects of the numerical representation within a token, attention
provides the mechanism through which information is exchanged across
different token positions. Understanding both operations is necessary
for reconstructing how a Transformer block processes a sequence.
